\reading{CR-6}{Credit Scoring and Rating}{STRIKE FIRST}{2}

\begin{crkeybox}[Learning objectives for this reading]
\begin{enumerate}[label=\textbf{\alph*.},leftmargin=2em]
\item Compare the credit scoring system to the credit rating system in assessing credit quality and describe the different types of each system.
\item Differentiate between through-the-cycle credit rating system and point-in-time credit rating system.
\item Describe the process for developing credit risk scoring and rating models.
\item Describe rating agencies' assignment methodologies for issue and issuer ratings, and identify the main criticisms of the credit rating agencies' ratings.
\end{enumerate}
\end{crkeybox}

% =====================================================================
\lo{CR-6 a}{Compare the credit scoring system to the credit rating system in assessing credit quality and describe the different types of each system.}

Both credit scoring and credit rating systems aim to quantify creditworthiness,
but they differ in their target audience, scale, and level of detail.

\begin{center}
\footnotesize
\begin{tabular}{L{2.6cm} L{5.8cm} L{6.2cm}}
\toprule
 & \thead{Credit scoring} & \thead{Credit rating} \\
\midrule
Who it covers & Smaller businesses and individuals &
  Larger businesses and governments \\
Scale & Numerical, typically 300 to 850. Higher numbers indicate a higher
  likelihood of repaying debt &
  Letter grades (e.g., AAA, AA, A, BBB, BBB+), with variations depending on the
  agency. AAA reflects the strongest repayment ability \\
Who produces it & Developed by banks \emph{internally}, analysing individual
  financial behaviours, payment history, outstanding debts, length of credit
  history, and types of credit used &
  Developed by \emph{external} credit rating agencies (e.g., Standard \& Poor's,
  Moody's, Fitch) based on comprehensive assessments of financial health,
  industry trends, and economic conditions \\
Update frequency & Updated more frequently & Updated less frequently \\
\bottomrule
\end{tabular}
\end{center}
\figcap{Scoring versus rating. Three things move together down the left column:
smaller borrowers, a numerical scale, and an internally produced, frequently
refreshed number.}

\begin{crkeybox}[The four benefits of a scoring or rating system]
\begin{enumerate}
\item Reduces subjectivity in loan evaluation and acceptance.
\item Enables analysis in credit risk management (e.g., scenario analysis and
stress testing).
\item Promotes transparency and consistency through a common framework.
\item Reduces the time for, and cost of, appraisal.
\end{enumerate}
\end{crkeybox}

% =====================================================================
\lo{CR-6 b}{Differentiate between through-the-cycle credit rating system and point-in-time credit rating system.}

\begin{center}
\footnotesize
\begin{tabular}{L{2.6cm} L{5.8cm} L{6.2cm}}
\toprule
 & \thead{Through-the-cycle (TTC)} & \thead{Point-in-time (PIT)} \\
\midrule
Used by & Major rating agencies & Banks, internally \\
Horizon & Takes a long-term view, over a full business cycle &
  Focuses on the current situation, possibly including the next year \\
Behaviour & Not updated frequently; less sensitive to short-term bumps, but may
  not reflect immediate risk &
  More volatile; reacts to short-term changes, potentially capturing default risk
  quickly \\
Better for & Long-term loans & Short-term loans \\
\bottomrule
\end{tabular}
\end{center}
\figcap{The pairing to hold: the system that is slower to move is the one used
externally by the agencies and suited to long-term loans; the system that reacts
fastest is the one banks run internally for short-term loans.}

% =====================================================================
\lo{CR-6 c}{Describe the process for developing credit risk scoring and rating models.}

Credit scores and ratings can be determined for either a specific \term{issuer}
(borrower) or a specific \term{issuance} (loan or bond). Issuers typically seek
credit scores or ratings when accessing capital markets to borrow money.

\subsection*{Scoring methods}

\begin{enumerate}
\item \term{Behavioral scoring.} Applies to consumer loans and observes the
historical financial behaviour of customers. Updates in real time based on new
behavioural information (e.g., payment behaviour, default estimates). Used for
credit limits, debt collection, and marketing new products.
\item \term{Profit scoring.} Focuses on the estimated profitability of a loan,
considering pricing, operational decisions, credit risk assessments, and default
predictions. Two types:
  \begin{itemize}
  \item \term{Account-level}: calculates profit for each individual account
  without considering connections between accounts.
  \item \term{Customer-level}: considers all accounts owned by a customer,
  allowing aggregation of profit across associated accounts.
  \end{itemize}
\end{enumerate}

\subsection*{Social lending (peer-to-peer lending)}

Enabled by advances in AI and fintech, allowing direct transactions between
lenders and borrowers without intermediaries. Credit decisions are made quickly,
but credit risk modelling in social lending is still developing.

Social lending often has higher default rates than traditional loans,
necessitating a better understanding of risks. Concerns include responsible use of
big data and transparency in this innovative lending channel.

\subsection*{The five-stage development process}

\begin{center}
\begin{tikzpicture}[font=\sffamily\footnotesize,
  stg/.style={draw=FRMNavy, fill=PaleNavy, line width=0.8pt, rounded corners=1pt,
              text width=4.0cm, align=center, minimum height=0.95cm}]
\node[stg] (s1) at (0,0)    {1. Data collection \& preprocessing};
\node[stg] (s2) at (0,-1.55) {2. Model fitting};
\node[stg] (s3) at (0,-3.10) {3. Validation};
\node[stg] (s4) at (0,-4.65) {4. Definition \& validation of risk ratings};
\node[stg] (s5) at (0,-6.30) {5. Implementation, monitoring, review};
\foreach \a/\b in {s1/s2, s2/s3, s3/s4}
  \draw[-{Stealth[length=5pt]}, FRMNavy, line width=0.9pt] (\a) -- (\b);
\draw[{Stealth[length=5pt]}-{Stealth[length=5pt]}, FRMNavy, line width=0.9pt] (s4) -- (s5);
% right-hand feedback: validation back to model fitting and to data collection
\draw[-{Stealth[length=4.5pt]}, FRMGold, line width=0.8pt]
  (s3.east) -- ++(1.5,0) |- (s2.east);
\draw[-{Stealth[length=4.5pt]}, FRMGold, line width=0.8pt]
  (s3.east) -- ++(2.4,0) |- (s1.east);
% left-hand feedback: implementation back up
\draw[-{Stealth[length=4.5pt]}, FRMGreen, line width=0.8pt]
  (s5.west) -- ++(-1.5,0) |- (s4.west);
\draw[-{Stealth[length=4.5pt]}, FRMGreen, line width=0.8pt]
  (s5.west) -- ++(-2.4,0) |- (s2.west);
\draw[-{Stealth[length=4.5pt]}, FRMGreen, line width=0.8pt]
  (s5.west) -- ++(-3.3,0) |- (s1.west);
\end{tikzpicture}
\end{center}
\figcap{The development process is a loop, not a pipeline. Two families of
feedback run backwards. The gold arrows on the right are validation failure:
stage 3 returns the process to model fitting and, if necessary, all the way to
data collection. The green arrows on the left are post-implementation
monitoring: degradation found at stage 5 can send the process back to any earlier
stage. Note also that stages 4 and 5 are joined by a two-way arrow, not a
one-way one.}

\subsection*{1) Data collection and preprocessing}

The first step is to collect the appropriate data and consider any data
preprocessing needs.

\begin{itemize}
\item \term{What data is needed.} Both successful loans and defaulted loans (over
90 days late) are included. Both quantitative (numbers) and qualitative
(descriptive) data are used. Subjectivity is minimised to create a reliable
scoring system.
\item \term{Data preprocessing.} This ensures data quality for analysis. Steps
include removing outliers (extreme values), selecting relevant factors (e.g.,
income, payment history), and transforming data for consistency (e.g., converting
currencies).
\end{itemize}

\begin{center}
\footnotesize
\begin{tabular}{L{6.9cm} L{7.6cm}}
\toprule
\thead{Data needed for consumer loans} & \thead{Data needed for corporate loans} \\
\midrule
Income; assets; existing debts; payment history; employment status; loan
collateral (e.g., car); loan type (e.g., mortgage) &
\textbf{Financial statements}, analysed to calculate ratios reflecting a
company's health: profitability ratios (e.g., return on equity) measure repayment
ability; solvency ratios (e.g., debt-to-equity) measure debt management;
liquidity ratios (e.g., current ratio) measure short-term obligation coverage;
management efficiency ratios assess operational performance.
\newline\newline
\textbf{Transaction data}, providing a more recent view of a borrower's financial
situation: reviews business transactions and payments; considers current debt,
delinquencies, and credit limits. \\
\bottomrule
\end{tabular}
\end{center}

\textbf{Other corporate considerations:}

\begin{enumerate}[label=\alph*)]
\item \term{Size and age.} Larger, established firms are seen as lower risk.
Market capitalization (public companies) or total assets (private companies) are
used as size indicators.
\item \term{Market conditions and competition.} Industry outlook, competition
level, and barriers to entry are assessed.
\item \term{Financial market data} (publicly traded companies only). Stock price,
volatility, and valuation data provide short-term insights.
\item \term{Corporate governance.} Board composition, executive qualifications
and accountability, ethical standards, and shareholder protection are reviewed.
\item \term{Corporate news and analytics} (emerging area). News articles, social
media, and other sources can offer valuable insights when used responsibly.
\end{enumerate}

\subsection*{2) Model fitting}

The second step is model fitting, which involves identifying model parameters
with the best descriptive ability (i.e., fit) for the model training data.

\begin{crfmlbox}[The fitted model and its optimisation]
Assume the following linear model:
\[ f(x) \;=\; \beta_0 + \beta_1 x_1 + \beta_2 x_2 + \dots + \beta_n x_n \]
The goal is to provide the best fit, that is, optimized estimates for $\beta$,
such that:
\[ \beta^{*} \;=\; \arg\min_{\beta \in \mathbb{R}} L(\beta, X) \]
\vk{The vector $\beta = (\beta_0, \beta_1, \dots, \beta_n)$ includes a constant
term $\beta_0$ and a series of coefficients $\beta_1,\dots,\beta_n$ for various
risk attributes. $\beta^{*}$ is the optimal parameter vector and $L$ is a loss
function representing the difference between the model's output and the
classification in the training data.}
\end{crfmlbox}

This is an empirical proxy for the true problem, which is to minimise default
losses. The formula is a regression-like process with a \emph{binary} (e.g.,
default / non-default) dependent variable rather than the continuous variable
normally used in standard regression models. Using algorithms, model fitting can
involve statistical tools, data mining, and operations research techniques.

\subsection*{3) Model validation}

Model validation checks the fitted model against a new sample, the validation
data, not included in the model training data. In other words, this is an
\term{out-of-sample test}. The validation data should share similar
characteristics (risk factors) with the training data, but involve different
borrowers or loans. If the validation results are not good, data collection,
preprocessing, or model fitting might need adjustments before re-validation.

Techniques like backtesting are used to compare the model's predictions with
actual outcomes. Backtesting can be done in two ways:

\begin{enumerate}[label=\alph*)]
\item \term{Out-of-sample}: uses recent data points not included in the initial
training data.
\item \term{Out-of-time}: uses data from a different time period than the
training data.
\end{enumerate}

\textbf{Enhancing validation:}

\begin{enumerate}
\item \term{Walk-forward testing.} Systematically tests the model on rolling time
periods. Rotate the evaluation time period forward systematically, incorporating
observations from previous periods (e.g., use data from periods $t-1$ and $t-2$
for testing period $t$).
\item \term{Resampling techniques.} Methods like bootstrapping and
cross-validation help create multiple training and validation sets to get a more
robust picture of the model's performance.
\item \term{Benchmarking.} Instead of just comparing predictions to actual
outcomes, benchmarking compares them to an external source such as credit ratings
from agencies. This approach does not aim to \emph{replicate} existing ratings
but to identify any discrepancies and understand why they occur. Deviations from
benchmarks might indicate areas for recalibration and model improvement.
\end{enumerate}

\subsection*{4) Definition and validation of ratings}

After data collection and preprocessing, model fitting, and model validation, the
derived credit scores should be mapped to a risk rating class.

\begin{itemize}
\item Each rating class should be associated with empirical probability of
default (PD) estimates.
\item Each rating class should be well diversified relative to its neighbouring
class, and the distribution of borrowers should not be highly concentrated in a
single risk grade.
\item The mapping process should be checked for consistency over time, given that
rating migration is a concern.
\item If this grouping does not produce adequate results, the process needs to be
re-evaluated and adjusted.
\end{itemize}

\subsection*{5) Implementation, monitoring, and review}

The final stage involves implementing the developed credit scoring or rating
system for real-time risk assessment of loan applications. Continuous monitoring,
review, and recalibration of the system are essential to identify any degradation
in risk estimates over time.

% =====================================================================
\lo{CR-6 d}{Describe rating agencies' assignment methodologies for issue and issuer ratings, and identify the main criticisms of the credit rating agencies' ratings.}

Credit scoring and rating systems are governed and imposed by regulators and
central banks. The Basel Committee has outlined the following five basic
specification requirements.

\begin{enumerate}
\item \term{Risk differentiation.} Grades must reflect actual risk differences,
not just meet capital requirements. Treat borrowers within the same grade
differently based on transaction-level variances. Ensure dispersion across grades
to avoid concentration risk.
\item \term{Continuous evaluation.} Borrower ratings should be regularly updated,
minimally on an annual basis, to reflect current risk profiles.
\item \term{Operational oversight.} Maintain system integrity through continuous
monitoring, efficient controls (e.g., stress testing), and feedback mechanisms
from stakeholders.
\item \term{Effective risk metrics.} Utilize risk factors that accurately predict
borrower creditworthiness, aiming for close alignment with actual outcomes.
\item \term{Data collection.} Incorporate comprehensive data representing
historical performance, past ratings, default estimates, and payment history to
ensure robust credit grading.
\end{enumerate}

\subsection*{Criticisms of credit rating agency scores}

\begin{enumerate}
\item \term{Lack of transparency.} Rating agency methodologies are often
proprietary, lacking public verification of fairness and validity.
\item \term{Potential conflicts of interest.} CRAs receive fees from the entities
they rate, raising concerns about impartiality.
\item \term{Promotion of debt expansion.} Enhanced risk management from CRAs may
reduce risk premiums, making debt more accessible, potentially fuelling credit
bubbles and crunches, and heightening systemic risk.
\item \term{Poor predictive ability.} CRAs have historically failed to predict
significant credit failures, such as Enron and Lehman Brothers, with studies
showing lower predictive ability compared to other methods.
\item \term{Procyclicality.} Despite claims of through-the-cycle modelling,
ratings tend to be overly optimistic during economic growth and overly
pessimistic during recessions, indicating some level of procyclicality. This does
not give a clear picture of a borrower's true risk.
\end{enumerate}

\begin{crtrapbox}[Procyclicality sits directly against CR-6 b]
The procyclicality criticism is aimed at the agencies, and the agencies are
precisely the users of \emph{through-the-cycle} ratings. The complaint is that
ratings move with the cycle despite the TTC label, not that they fail to move.
\end{crtrapbox}
